
\documentclass{article}
\usepackage{colortbl}
\usepackage{makecell}
\usepackage{multirow}
\usepackage{supertabular}

\begin{document}

\newcounter{utterance}

\centering \large Interaction Transcript for game `matchit\_ascii', experiment `similar\_grid\_1', episode 4 with glm{-}5.2.
\vspace{24pt}

{ \footnotesize  \setcounter{utterance}{1}
\setlength{\tabcolsep}{0pt}
\begin{supertabular}{c@{$\;$}|p{.15\linewidth}@{}p{.15\linewidth}p{.15\linewidth}p{.15\linewidth}p{.15\linewidth}p{.15\linewidth}}
    \# & $\;$A & \multicolumn{4}{c}{Game Master} & $\;\:$B\\
    \hline

    \theutterance \stepcounter{utterance}  
    & & \multicolumn{4}{p{0.6\linewidth}}{
        \cellcolor[rgb]{0.9,0.9,0.9}{
            \makecell[{{p{\linewidth}}}]{
                \texttt{\tiny{[P1$\langle$GM]}}
                \texttt{Stai partecipando a un gioco di indovinelli collaborativo.} \\
\texttt{L'obiettivo è scoprire se questa griglia e un'altra immagine che vedo solo io sono uguali.} \\
\texttt{Per favore descrivi prima la tua griglia. Poi, io fornirò la mia descrizione e potremo farci domande reciproche sulle griglie per capire se sono uguali. Ora inizia la tua breve descrizione della griglia con "DESCRIZIONE:" seguito dalla descrizione. Non aggiungere altro. Questa è la tua griglia:} \\
\texttt{X X ▢ X X} \\
\texttt{X ▢ ▢ ▢ X} \\
\texttt{▢ ▢ ▢ ▢ ▢} \\
\texttt{X X ▢ X X} \\
\texttt{X X ▢ X X} \\
            }
        }
    }
    & \\ \\

    \theutterance \stepcounter{utterance}  
    & \multicolumn{4}{p{0.6\linewidth}}{
        \cellcolor[rgb]{0.9,0.9,0.9}{
            \makecell[{{p{\linewidth}}}]{
                \texttt{\tiny{[P1$\rangle$GM]}}
                \texttt{DESCRIZIONE: Griglia 5x5 simmetrica. Cornice esterna di "X", tranne al centro dei lati dove c'è "▢". Croce centrale di "▢" nella riga 3 e colonna 3, tutto il resto è "X".} \\
            }
        }
    }
    & & \\ \\

    \theutterance \stepcounter{utterance}  
    & & & \multicolumn{2}{p{0.3\linewidth}}{
        \cellcolor[rgb]{0.9,0.9,0.9}{
            \makecell[{{p{\linewidth}}}]{
                \texttt{\tiny{[GM$|$GM]}}
                \texttt{continue} \\
            }
        }
    }
    & & \\ \\

    \theutterance \stepcounter{utterance}  
    & & \multicolumn{4}{p{0.6\linewidth}}{
        \cellcolor[rgb]{0.9,0.9,0.9}{
            \makecell[{{p{\linewidth}}}]{
                \texttt{\tiny{[GM$\rangle$P2]}}
                \texttt{Stai partecipando a un gioco di indovinelli collaborativo.} \\
\texttt{L'obiettivo è scoprire se questa griglia e un'altra immagine che vedo solo io sono uguali.} \\
\texttt{Per favore descrivi prima la tua griglia. Poi, io fornirò la mia descrizione e potremo farci domande reciproche sulle griglie per capire se sono uguali. Ora inizia la tua breve descrizione della griglia con "DESCRIZIONE:" seguito dalla descrizione. Non aggiungere altro. Questa è la tua griglia:} \\
\texttt{X X ▢ X X} \\
\texttt{X ▢ ▢ X X} \\
\texttt{▢ ▢ ▢ ▢ ▢} \\
\texttt{X ▢ ▢ X X} \\
\texttt{X X ▢ X X} \\
            }
        }
    }
    & \\ \\

    \theutterance \stepcounter{utterance}  
    & & & \multicolumn{4}{p{0.6\linewidth}}{
        \cellcolor[rgb]{0.9,0.9,0.9}{
            \makecell[{{p{\linewidth}}}]{
                \texttt{\tiny{[GM$\langle$P2]}}
                \texttt{DESCRIZIONE: Una griglia 5x5 simmetrica sia orizzontalmente che verticalmente. La terza riga e la terza colonna sono formate interamente da quadrati vuoti (▢). I quadrati rimanenti, che formano un bordo spesso due, contengono la lettera X.} \\
            }
        }
    }
     \\ \\

    \theutterance \stepcounter{utterance}  
    & & & \multicolumn{2}{p{0.3\linewidth}}{
        \cellcolor[rgb]{0.9,0.9,0.9}{
            \makecell[{{p{\linewidth}}}]{
                \texttt{\tiny{[GM$|$GM]}}
                \texttt{continue} \\
            }
        }
    }
    & & \\ \\

    \theutterance \stepcounter{utterance}  
    & & \multicolumn{4}{p{0.6\linewidth}}{
        \cellcolor[rgb]{0.9,0.9,0.9}{
            \makecell[{{p{\linewidth}}}]{
                \texttt{\tiny{[GM$\rangle$P2]}}
                \texttt{Questa è la miaDESCRIZIONE: Griglia 5x5 simmetrica. Cornice esterna di "X", tranne al centro dei lati dove c'è "▢". Croce centrale di "▢" nella riga 3 e colonna 3, tutto il resto è "X".} \\
\texttt{Ora fai una domanda per scoprire nuovi aspetti della mia griglia che potrebbero essere diversi dalla tua griglia. Inizia con "DOMANDA:" e non aggiungere altro.} \\
            }
        }
    }
    & \\ \\

    \theutterance \stepcounter{utterance}  
    & & & \multicolumn{4}{p{0.6\linewidth}}{
        \cellcolor[rgb]{0.9,0.9,0.9}{
            \makecell[{{p{\linewidth}}}]{
                \texttt{\tiny{[GM$\langle$P2]}}
                \texttt{DOMANDA: Nella tua descrizione menzioni una "cornice esterna" di X con eccezioni al centro dei lati. Se consideriamo la riga 1 (che per me è X X ▢ X X), anche le celle in posizione 1 e 5 (i "corni" della cornice) sono X? Inoltre, le celle che non fanno parte della croce centrale e non sono sul bordo esterno (come ad esempio la cella in riga 2, colonna 2) sono X?} \\
            }
        }
    }
     \\ \\

    \theutterance \stepcounter{utterance}  
    & & & \multicolumn{2}{p{0.3\linewidth}}{
        \cellcolor[rgb]{0.9,0.9,0.9}{
            \makecell[{{p{\linewidth}}}]{
                \texttt{\tiny{[GM$|$GM]}}
                \texttt{continue} \\
            }
        }
    }
    & & \\ \\

    \theutterance \stepcounter{utterance}  
    & & \multicolumn{4}{p{0.6\linewidth}}{
        \cellcolor[rgb]{0.9,0.9,0.9}{
            \makecell[{{p{\linewidth}}}]{
                \texttt{\tiny{[P1$\langle$GM]}}
                \texttt{Questa è la miaDESCRIZIONE: Una griglia 5x5 simmetrica sia orizzontalmente che verticalmente. La terza riga e la terza colonna sono formate interamente da quadrati vuoti (▢). I quadrati rimanenti, che formano un bordo spesso due, contengono la lettera X.} \\
\texttt{DOMANDA: Nella tua descrizione menzioni una "cornice esterna" di X con eccezioni al centro dei lati. Se consideriamo la riga 1 (che per me è X X ▢ X X), anche le celle in posizione 1 e 5 (i "corni" della cornice) sono X? Inoltre, le celle che non fanno parte della croce centrale e non sono sul bordo esterno (come ad esempio la cella in riga 2, colonna 2) sono X?Inizia la tua risposta con "RISPOSTA:" e non aggiungere altro.} \\
            }
        }
    }
    & \\ \\

    \theutterance \stepcounter{utterance}  
    & \multicolumn{4}{p{0.6\linewidth}}{
        \cellcolor[rgb]{0.9,0.9,0.9}{
            \makecell[{{p{\linewidth}}}]{
                \texttt{\tiny{[P1$\rangle$GM]}}
                \texttt{RISPOSTA: Sì, le celle in posizione 1 e 5 della riga 1 sono X. Sì, le celle interne non appartenenti alla croce centrale (come riga 2, colonna 2) sono X.} \\
            }
        }
    }
    & & \\ \\

    \theutterance \stepcounter{utterance}  
    & & & \multicolumn{2}{p{0.3\linewidth}}{
        \cellcolor[rgb]{0.9,0.9,0.9}{
            \makecell[{{p{\linewidth}}}]{
                \texttt{\tiny{[GM$|$GM]}}
                \texttt{continue} \\
            }
        }
    }
    & & \\ \\

    \theutterance \stepcounter{utterance}  
    & & \multicolumn{4}{p{0.6\linewidth}}{
        \cellcolor[rgb]{0.9,0.9,0.9}{
            \makecell[{{p{\linewidth}}}]{
                \texttt{\tiny{[P1$\langle$GM]}}
                \texttt{Ora fai una domanda per scoprire nuovi aspetti della mia griglia che potrebbero essere diversi dalla tua griglia. Inizia con "DOMANDA:" e non aggiungere altro.} \\
            }
        }
    }
    & \\ \\

    \theutterance \stepcounter{utterance}  
    & \multicolumn{4}{p{0.6\linewidth}}{
        \cellcolor[rgb]{0.9,0.9,0.9}{
            \makecell[{{p{\linewidth}}}]{
                \texttt{\tiny{[P1$\rangle$GM]}}
                \texttt{DOMANDA: La cella nell'angolo in alto a sinistra (riga 1, colonna 1) e quella nell'angolo in basso a destra (riga 5, colonna 5) contengono la lettera X?} \\
            }
        }
    }
    & & \\ \\

    \theutterance \stepcounter{utterance}  
    & & & \multicolumn{2}{p{0.3\linewidth}}{
        \cellcolor[rgb]{0.9,0.9,0.9}{
            \makecell[{{p{\linewidth}}}]{
                \texttt{\tiny{[GM$|$GM]}}
                \texttt{continue} \\
            }
        }
    }
    & & \\ \\

    \theutterance \stepcounter{utterance}  
    & & \multicolumn{4}{p{0.6\linewidth}}{
        \cellcolor[rgb]{0.9,0.9,0.9}{
            \makecell[{{p{\linewidth}}}]{
                \texttt{\tiny{[GM$\rangle$P2]}}
                \texttt{RISPOSTA: Sì, le celle in posizione 1 e 5 della riga 1 sono X. Sì, le celle interne non appartenenti alla croce centrale (come riga 2, colonna 2) sono X.} \\
\texttt{DOMANDA: La cella nell'angolo in alto a sinistra (riga 1, colonna 1) e quella nell'angolo in basso a destra (riga 5, colonna 5) contengono la lettera X?Inizia la tua risposta con "RISPOSTA:" e non aggiungere altro.} \\
            }
        }
    }
    & \\ \\

    \theutterance \stepcounter{utterance}  
    & & & \multicolumn{4}{p{0.6\linewidth}}{
        \cellcolor[rgb]{0.9,0.9,0.9}{
            \makecell[{{p{\linewidth}}}]{
                \texttt{\tiny{[GM$\langle$P2]}}
                \texttt{RISPOSTA: Sì, entrambe le celle nell'angolo in alto a sinistra (riga 1, colonna 1) e nell'angolo in basso a destra (riga 5, colonna 5) contengono la lettera X.} \\
            }
        }
    }
     \\ \\

    \theutterance \stepcounter{utterance}  
    & & & \multicolumn{2}{p{0.3\linewidth}}{
        \cellcolor[rgb]{0.9,0.9,0.9}{
            \makecell[{{p{\linewidth}}}]{
                \texttt{\tiny{[GM$|$GM]}}
                \texttt{continue} \\
            }
        }
    }
    & & \\ \\

    \theutterance \stepcounter{utterance}  
    & & \multicolumn{4}{p{0.6\linewidth}}{
        \cellcolor[rgb]{0.9,0.9,0.9}{
            \makecell[{{p{\linewidth}}}]{
                \texttt{\tiny{[GM$\rangle$P2]}}
                \texttt{Ora fai una domanda per scoprire nuovi aspetti della mia griglia che potrebbero essere diversi dalla tua griglia. Inizia con "DOMANDA:" e non aggiungere altro.} \\
            }
        }
    }
    & \\ \\

    \theutterance \stepcounter{utterance}  
    & & & \multicolumn{4}{p{0.6\linewidth}}{
        \cellcolor[rgb]{0.9,0.9,0.9}{
            \makecell[{{p{\linewidth}}}]{
                \texttt{\tiny{[GM$\langle$P2]}}
                \texttt{DOMANDA: Sulla base delle tue risposte precedenti sembra che la nostra griglia sia identica. Per confermare in modo definitivo, puoi elencare il contenuto esatto della riga 2 e della riga 4 della tua griglia?} \\
            }
        }
    }
     \\ \\

    \theutterance \stepcounter{utterance}  
    & & & \multicolumn{2}{p{0.3\linewidth}}{
        \cellcolor[rgb]{0.9,0.9,0.9}{
            \makecell[{{p{\linewidth}}}]{
                \texttt{\tiny{[GM$|$GM]}}
                \texttt{continue} \\
            }
        }
    }
    & & \\ \\

    \theutterance \stepcounter{utterance}  
    & & \multicolumn{4}{p{0.6\linewidth}}{
        \cellcolor[rgb]{0.9,0.9,0.9}{
            \makecell[{{p{\linewidth}}}]{
                \texttt{\tiny{[P1$\langle$GM]}}
                \texttt{RISPOSTA: Sì, entrambe le celle nell'angolo in alto a sinistra (riga 1, colonna 1) e nell'angolo in basso a destra (riga 5, colonna 5) contengono la lettera X.} \\
\texttt{DOMANDA: Sulla base delle tue risposte precedenti sembra che la nostra griglia sia identica. Per confermare in modo definitivo, puoi elencare il contenuto esatto della riga 2 e della riga 4 della tua griglia?Inizia la tua risposta con "RISPOSTA:" e non aggiungere altro.} \\
            }
        }
    }
    & \\ \\

    \theutterance \stepcounter{utterance}  
    & \multicolumn{4}{p{0.6\linewidth}}{
        \cellcolor[rgb]{0.9,0.9,0.9}{
            \makecell[{{p{\linewidth}}}]{
                \texttt{\tiny{[P1$\rangle$GM]}}
                \texttt{RISPOSTA: Riga 2: X ▢ ▢ ▢ X. Riga 4: X X ▢ X X.} \\
            }
        }
    }
    & & \\ \\

    \theutterance \stepcounter{utterance}  
    & & & \multicolumn{2}{p{0.3\linewidth}}{
        \cellcolor[rgb]{0.9,0.9,0.9}{
            \makecell[{{p{\linewidth}}}]{
                \texttt{\tiny{[GM$|$GM]}}
                \texttt{continue} \\
            }
        }
    }
    & & \\ \\

    \theutterance \stepcounter{utterance}  
    & & \multicolumn{4}{p{0.6\linewidth}}{
        \cellcolor[rgb]{0.9,0.9,0.9}{
            \makecell[{{p{\linewidth}}}]{
                \texttt{\tiny{[P1$\langle$GM]}}
                \texttt{Ora fai una domanda per scoprire nuovi aspetti della mia griglia che potrebbero essere diversi dalla tua griglia. Inizia con "DOMANDA:" e non aggiungere altro.} \\
            }
        }
    }
    & \\ \\

    \theutterance \stepcounter{utterance}  
    & \multicolumn{4}{p{0.6\linewidth}}{
        \cellcolor[rgb]{0.9,0.9,0.9}{
            \makecell[{{p{\linewidth}}}]{
                \texttt{\tiny{[P1$\rangle$GM]}}
                \texttt{DOMANDA: Per un'ultima verifica, puoi elencare il contenuto esatto della riga 1 e della riga 3 della tua griglia?} \\
            }
        }
    }
    & & \\ \\

    \theutterance \stepcounter{utterance}  
    & & & \multicolumn{2}{p{0.3\linewidth}}{
        \cellcolor[rgb]{0.9,0.9,0.9}{
            \makecell[{{p{\linewidth}}}]{
                \texttt{\tiny{[GM$|$GM]}}
                \texttt{continue} \\
            }
        }
    }
    & & \\ \\

    \theutterance \stepcounter{utterance}  
    & & \multicolumn{4}{p{0.6\linewidth}}{
        \cellcolor[rgb]{0.9,0.9,0.9}{
            \makecell[{{p{\linewidth}}}]{
                \texttt{\tiny{[GM$\rangle$P2]}}
                \texttt{RISPOSTA: Riga 2: X ▢ ▢ ▢ X. Riga 4: X X ▢ X X.} \\
\texttt{DOMANDA: Per un'ultima verifica, puoi elencare il contenuto esatto della riga 1 e della riga 3 della tua griglia?Inizia la tua risposta con "RISPOSTA:" e non aggiungere altro.} \\
            }
        }
    }
    & \\ \\

    \theutterance \stepcounter{utterance}  
    & & & \multicolumn{4}{p{0.6\linewidth}}{
        \cellcolor[rgb]{0.9,0.9,0.9}{
            \makecell[{{p{\linewidth}}}]{
                \texttt{\tiny{[GM$\langle$P2]}}
                \texttt{RISPOSTA: Riga 1: X X ▢ X X. Riga 3: ▢ ▢ ▢ ▢ ▢.} \\
            }
        }
    }
     \\ \\

    \theutterance \stepcounter{utterance}  
    & & & \multicolumn{2}{p{0.3\linewidth}}{
        \cellcolor[rgb]{0.9,0.9,0.9}{
            \makecell[{{p{\linewidth}}}]{
                \texttt{\tiny{[GM$|$GM]}}
                \texttt{continue} \\
            }
        }
    }
    & & \\ \\

    \theutterance \stepcounter{utterance}  
    & & \multicolumn{4}{p{0.6\linewidth}}{
        \cellcolor[rgb]{0.9,0.9,0.9}{
            \makecell[{{p{\linewidth}}}]{
                \texttt{\tiny{[GM$\rangle$P2]}}
                \texttt{Ora fai una domanda per scoprire nuovi aspetti della mia griglia che potrebbero essere diversi dalla tua griglia. Inizia con "DOMANDA:" e non aggiungere altro.} \\
            }
        }
    }
    & \\ \\

    \theutterance \stepcounter{utterance}  
    & & & \multicolumn{4}{p{0.6\linewidth}}{
        \cellcolor[rgb]{0.9,0.9,0.9}{
            \makecell[{{p{\linewidth}}}]{
                \texttt{\tiny{[GM$\langle$P2]}}
                \texttt{DOMANDA: La tua risposta precedente indica che la riga 4 è "X X ▢ X X", mentre la riga 2 è "X ▢ ▢ ▢ X". Nella mia griglia, sia la riga 2 che la riga 4 sono "X ▢ ▢ X X". Puoi confermare se la riga 4 della tua griglia è effettivamente "X X ▢ X X" e se la cella in riga 4, colonna 2 è una "X" o un "▢"?} \\
            }
        }
    }
     \\ \\

    \theutterance \stepcounter{utterance}  
    & & & \multicolumn{2}{p{0.3\linewidth}}{
        \cellcolor[rgb]{0.9,0.9,0.9}{
            \makecell[{{p{\linewidth}}}]{
                \texttt{\tiny{[GM$|$GM]}}
                \texttt{continue} \\
            }
        }
    }
    & & \\ \\

    \theutterance \stepcounter{utterance}  
    & & \multicolumn{4}{p{0.6\linewidth}}{
        \cellcolor[rgb]{0.9,0.9,0.9}{
            \makecell[{{p{\linewidth}}}]{
                \texttt{\tiny{[P1$\langle$GM]}}
                \texttt{RISPOSTA: Riga 1: X X ▢ X X. Riga 3: ▢ ▢ ▢ ▢ ▢.} \\
\texttt{DOMANDA: La tua risposta precedente indica che la riga 4 è "X X ▢ X X", mentre la riga 2 è "X ▢ ▢ ▢ X". Nella mia griglia, sia la riga 2 che la riga 4 sono "X ▢ ▢ X X". Puoi confermare se la riga 4 della tua griglia è effettivamente "X X ▢ X X" e se la cella in riga 4, colonna 2 è una "X" o un "▢"?Inizia la tua risposta con "RISPOSTA:" e non aggiungere altro.} \\
            }
        }
    }
    & \\ \\

    \theutterance \stepcounter{utterance}  
    & \multicolumn{4}{p{0.6\linewidth}}{
        \cellcolor[rgb]{0.9,0.9,0.9}{
            \makecell[{{p{\linewidth}}}]{
                \texttt{\tiny{[P1$\rangle$GM]}}
                \texttt{RISPOSTA: Confermo che la riga 4 della mia griglia è "X X ▢ X X". La cella in riga 4, colonna 2 è una "X".} \\
            }
        }
    }
    & & \\ \\

    \theutterance \stepcounter{utterance}  
    & & & \multicolumn{2}{p{0.3\linewidth}}{
        \cellcolor[rgb]{0.9,0.9,0.9}{
            \makecell[{{p{\linewidth}}}]{
                \texttt{\tiny{[GM$|$GM]}}
                \texttt{continue} \\
            }
        }
    }
    & & \\ \\

    \theutterance \stepcounter{utterance}  
    & & \multicolumn{4}{p{0.6\linewidth}}{
        \cellcolor[rgb]{0.9,0.9,0.9}{
            \makecell[{{p{\linewidth}}}]{
                \texttt{\tiny{[P1$\langle$GM]}}
                \texttt{Ora fai una domanda per scoprire nuovi aspetti della mia griglia che potrebbero essere diversi dalla tua griglia. Inizia con "DOMANDA:" e non aggiungere altro.} \\
            }
        }
    }
    & \\ \\

    \theutterance \stepcounter{utterance}  
    & \multicolumn{4}{p{0.6\linewidth}}{
        \cellcolor[rgb]{0.9,0.9,0.9}{
            \makecell[{{p{\linewidth}}}]{
                \texttt{\tiny{[P1$\rangle$GM]}}
                \texttt{DOMANDA: Nella tua affermazione precedente hai indicato che la tua riga 2 e la riga 4 sono "X ▢ ▢ X X", il che differisce dalla mia riga 2 ("X ▢ ▢ ▢ X"). Puoi confermare l'esatto contenuto della tua riga 2 e specificamente se la cella in riga 2, colonna 4 è una "X" o un "▢"?} \\
            }
        }
    }
    & & \\ \\

    \theutterance \stepcounter{utterance}  
    & & & \multicolumn{2}{p{0.3\linewidth}}{
        \cellcolor[rgb]{0.9,0.9,0.9}{
            \makecell[{{p{\linewidth}}}]{
                \texttt{\tiny{[GM$|$GM]}}
                \texttt{continue} \\
            }
        }
    }
    & & \\ \\

    \theutterance \stepcounter{utterance}  
    & & \multicolumn{4}{p{0.6\linewidth}}{
        \cellcolor[rgb]{0.9,0.9,0.9}{
            \makecell[{{p{\linewidth}}}]{
                \texttt{\tiny{[GM$\rangle$P2]}}
                \texttt{RISPOSTA: Confermo che la riga 4 della mia griglia è "X X ▢ X X". La cella in riga 4, colonna 2 è una "X".} \\
\texttt{DOMANDA: Nella tua affermazione precedente hai indicato che la tua riga 2 e la riga 4 sono "X ▢ ▢ X X", il che differisce dalla mia riga 2 ("X ▢ ▢ ▢ X"). Puoi confermare l'esatto contenuto della tua riga 2 e specificamente se la cella in riga 2, colonna 4 è una "X" o un "▢"?Inizia la tua risposta con "RISPOSTA:" e non aggiungere altro.} \\
            }
        }
    }
    & \\ \\

    \theutterance \stepcounter{utterance}  
    & & & \multicolumn{4}{p{0.6\linewidth}}{
        \cellcolor[rgb]{0.9,0.9,0.9}{
            \makecell[{{p{\linewidth}}}]{
                \texttt{\tiny{[GM$\langle$P2]}}
                \texttt{RISPOSTA: La mia riga 2 è "X ▢ ▢ X X", quindi la cella in riga 2, colonna 4 è una "X".} \\
            }
        }
    }
     \\ \\

    \theutterance \stepcounter{utterance}  
    & & & \multicolumn{2}{p{0.3\linewidth}}{
        \cellcolor[rgb]{0.9,0.9,0.9}{
            \makecell[{{p{\linewidth}}}]{
                \texttt{\tiny{[GM$|$GM]}}
                \texttt{continue} \\
            }
        }
    }
    & & \\ \\

    \theutterance \stepcounter{utterance}  
    & & \multicolumn{4}{p{0.6\linewidth}}{
        \cellcolor[rgb]{0.9,0.9,0.9}{
            \makecell[{{p{\linewidth}}}]{
                \texttt{\tiny{[GM$\rangle$P2]}}
                \texttt{Ora prendi una decisione. Cosa ne pensi: la tua griglia e l'altra griglia descritta sono uguali?  Scrivi "DECISIONE: stessa griglia" se pensi che siano la stessa griglia oppure "DECISIONE: griglie diverse" se pensi che siano griglie diverse. Non aggiungere nient'altro.} \\
            }
        }
    }
    & \\ \\

    \theutterance \stepcounter{utterance}  
    & & & \multicolumn{4}{p{0.6\linewidth}}{
        \cellcolor[rgb]{0.9,0.9,0.9}{
            \makecell[{{p{\linewidth}}}]{
                \texttt{\tiny{[GM$\langle$P2]}}
                \texttt{DECISIONE: griglie diverse} \\
            }
        }
    }
     \\ \\

    \theutterance \stepcounter{utterance}  
    & & & \multicolumn{2}{p{0.3\linewidth}}{
        \cellcolor[rgb]{0.9,0.9,0.9}{
            \makecell[{{p{\linewidth}}}]{
                \texttt{\tiny{[GM$|$GM]}}
                \texttt{continue} \\
            }
        }
    }
    & & \\ \\

    \theutterance \stepcounter{utterance}  
    & & & \multicolumn{2}{p{0.3\linewidth}}{
        \cellcolor[rgb]{0.9,0.9,0.9}{
            \makecell[{{p{\linewidth}}}]{
                \texttt{\tiny{[GM$|$GM]}}
                \texttt{success} \\
            }
        }
    }
    & & \\ \\

    \theutterance \stepcounter{utterance}  
    & & \multicolumn{4}{p{0.6\linewidth}}{
        \cellcolor[rgb]{0.9,0.9,0.9}{
            \makecell[{{p{\linewidth}}}]{
                \texttt{\tiny{[P1$\langle$GM]}}
                \texttt{RISPOSTA: La mia riga 2 è "X ▢ ▢ X X", quindi la cella in riga 2, colonna 4 è una "X".} \\
\texttt{Ora prendi una decisione. Cosa ne pensi: la tua griglia e l'altra griglia descritta sono uguali?  Scrivi "DECISIONE: stessa griglia" se pensi che siano la stessa griglia oppure "DECISIONE: griglie diverse" se pensi che siano griglie diverse. Non aggiungere nient'altro.} \\
            }
        }
    }
    & \\ \\

    \theutterance \stepcounter{utterance}  
    & \multicolumn{4}{p{0.6\linewidth}}{
        \cellcolor[rgb]{0.9,0.9,0.9}{
            \makecell[{{p{\linewidth}}}]{
                \texttt{\tiny{[P1$\rangle$GM]}}
                \texttt{DECISIONE: griglie diverse} \\
            }
        }
    }
    & & \\ \\

    \theutterance \stepcounter{utterance}  
    & & & \multicolumn{2}{p{0.3\linewidth}}{
        \cellcolor[rgb]{0.9,0.9,0.9}{
            \makecell[{{p{\linewidth}}}]{
                \texttt{\tiny{[GM$|$GM]}}
                \texttt{continue} \\
            }
        }
    }
    & & \\ \\

    \theutterance \stepcounter{utterance}  
    & & & \multicolumn{2}{p{0.3\linewidth}}{
        \cellcolor[rgb]{0.9,0.9,0.9}{
            \makecell[{{p{\linewidth}}}]{
                \texttt{\tiny{[GM$|$GM]}}
                \texttt{success} \\
            }
        }
    }
    & & \\ \\

\end{supertabular}
}

\end{document}
